\documentclass[10pt,a4paper]{amsart}
\usepackage{amsmath,amssymb,amsthm,graphicx}
\usepackage{verbatim}
\newenvironment{reference}{\comment}{\endcomment}
%\newenvironment{reference}{}{}
\newenvironment{slogan}{\comment}{\endcomment}
\newenvironment{history}{\comment}{\endcomment}

% For commutative diagrams we use Xy-pic
\usepackage[all]{xy}

% We use 2cell for 2-commutative diagrams.
\xyoption{2cell}
% Theorem environments.
%
\theoremstyle{plain}
\newtheorem{theorem}[subsection]{Theorem}
\newtheorem{proposition}[subsection]{Proposition}
\newtheorem{lemma}[subsection]{Lemma}

\theoremstyle{definition}
\newtheorem{definition}[subsection]{Definition}
\newtheorem{example}[subsection]{Example}
\newtheorem{exercise}[subsection]{Exercise}
\newtheorem{situation}[subsection]{Situation}

\theoremstyle{remark}
\newtheorem{remark}[subsection]{Remark}
\newtheorem{remarks}[subsection]{Remarks}

\numberwithin{equation}{subsection}

% Macros
%
\def\lim{\mathop{\mathrm{lim}}\nolimits}
\def\colim{\mathop{\mathrm{colim}}\nolimits}
\def\Spec{\mathop{\mathrm{Spec}}}
\def\Hom{\mathop{\mathrm{Hom}}\nolimits}
\def\Ext{\mathop{\mathrm{Ext}}\nolimits}
\def\SheafHom{\mathop{\mathcal{H}\!\mathit{om}}\nolimits}
\def\SheafExt{\mathop{\mathcal{E}\!\mathit{xt}}\nolimits}
\def\Sch{\mathit{Sch}}
\def\Mor{\mathop{\mathrm{Mor}}\nolimits}
\def\Ob{\mathop{\mathrm{Ob}}\nolimits}
\def\Sh{\mathop{\mathit{Sh}}\nolimits}
\def\NL{\mathop{N\!L}\nolimits}
\def\CH{\mathop{\mathrm{CH}}\nolimits}
\def\proetale{{pro\text{-}\acute{e}tale}}
\def\etale{{\acute{e}tale}}
\def\QCoh{\mathit{QCoh}}
\def\Ker{\mathop{\mathrm{Ker}}}
\def\Im{\mathop{\mathrm{Im}}}
\def\Coker{\mathop{\mathrm{Coker}}}
\def\Coim{\mathop{\mathrm{Coim}}}

% Boxtimes
%
\DeclareMathSymbol{\boxtimes}{\mathbin}{AMSa}{"02}

%
% Macros for moduli stacks/spaces
%
\def\QCohstack{\mathcal{QC}\!\mathit{oh}}
\def\Cohstack{\mathcal{C}\!\mathit{oh}}
\def\Spacesstack{\mathcal{S}\!\mathit{paces}}
\def\Quotfunctor{\mathrm{Quot}}
\def\Hilbfunctor{\mathrm{Hilb}}
\def\Curvesstack{\mathcal{C}\!\mathit{urves}}
\def\Polarizedstack{\mathcal{P}\!\mathit{olarized}}
\def\Complexesstack{\mathcal{C}\!\mathit{omplexes}}
% \Pic is the operator that assigns to X its picard group, usage \Pic(X)
% \Picardstack_{X/B} denotes the Picard stack of X over B
% \Picardfunctor_{X/B} denotes the Picard functor of X over B
\def\Pic{\mathop{\mathrm{Pic}}\nolimits}
\def\Picardstack{\mathcal{P}\!\mathit{ic}}
\def\Picardfunctor{\mathrm{Pic}}
\def\Deformationcategory{\mathcal{D}\!\mathit{ef}}
\usepackage{cleveref}
\crefname{originaldefinition}{original definition}{original definitions}
\Crefname{originaldefinition}{Original definition}{Original definitions}
\makeatletter
\newlabel{definition-cohomology-shift}{{011H}{source}{}{sourcepointer-1}{}}
\newlabel{definition-cohomology-shift@cref}{{[originaldefinition][0][]011H}{[1][0][]source}{}{}{}}
\makeatother
\crefname{originallemma}{original lemma}{original lemmas}
\Crefname{originallemma}{Original lemma}{Original lemmas}
\makeatletter
\newlabel{lemma-snake}{{010H}{source}{}{sourcepointer-2}{}}
\newlabel{lemma-snake@cref}{{[originallemma][0][]010H}{[1][0][]source}{}{}{}}
\makeatother
\crefname{originallemma}{original lemma}{original lemmas}
\Crefname{originallemma}{Original lemma}{Original lemmas}
\makeatletter
\newlabel{lemma-snake-natural}{{08N7}{source}{}{sourcepointer-3}{}}
\newlabel{lemma-snake-natural@cref}{{[originallemma][0][]08N7}{[1][0][]source}{}{}{}}
\makeatother
\title{The long exact cohomology sequence of a short exact sequence of cochain complexes}
\author{The Stacks Project Authors}
\date{stacks-project:homology.tex:local-snapshot-20261012; GNU FDL 1.2 or later}
\begin{document}
\maketitle
\noindent\textit{Editorial note.} All delivered mathematical text is the author's own, copied from the pinned Stacks Project chapter file homology.tex; no statement or proof step was rewritten, reordered or removed.\par\medskip
\noindent\textit{Editorial note.} The proof is the author's complete argument; the diagram to which the snake lemma is applied is displayed in the counted span, the identifications of its outer terms are stated, and the closing remarks about functoriality and shifts are conclusions rather than deferrals.\par\medskip
\noindent\textit{Editorial note.} Every labelled reference occurring inside the counted spans is carried in assembly.external\_references as a non-load-bearing navigation pointer (same-chapter ones at their definition site, cross-chapter ones at the invocation inside the counted span), so no load-bearing mathematics is replaced by a stub. No bibliography citation occurs inside the counted spans.\par\medskip
\noindent\textit{Editorial note.} The same-name check for this card scanned every staged Stacks chapter of this repository -- all sources/stacks-*/pinned-original-source.tex files, the card's own chapter excluded -- and the labels found are recorded in dedup.same\_name\_cross\_chapter on the card.\par\medskip
\noindent\textit{Editorial note.} The preamble is the chapter's own native preamble excerpt already used for this source: the comment-type reference, slogan and history environments (native preamble lines 8-13), the Xy-pic setup (native preamble lines 14-18) and the native theorem declarations with the mathematical macros (native preamble lines 154-217, the actual last line of the file).\par\medskip
\noindent\textit{Allowed standard prerequisites (editorial).}\begin{itemize}
\item The snake lemma. Same chapter, Lemma lemma-snake. Supplies the connecting map of the displayed diagram together with the exactness of the resulting six-term sequence, which is exactly the exactness claimed for the long sequence.
\item Naturality of the connecting map. Same chapter, Lemma lemma-snake-natural. Supplies the functoriality of the constructed boundary map in the data of the short exact sequence.
\item Cohomology and its shift. Same chapter, the definition of the cohomology groups of a cochain complex and of the shift in degree. Supply the objects of the displayed sequence and the meaning of compatibility with shifts.
\end{itemize}\medskip
\ifdefined\hypertarget\hypertarget{sourcepointer-1}{}\fi
\noindent\textit{Source reference.} Original definition 011H: sources/stacks-homology/pinned-original-source.tex, source lines 3679--3679. Navigation pointer invoked at proof line 3437; the counted span contains the pointer only and the cited result is not reproduced inside it, so the pointer is registered as a non-load-bearing navigation stub in the established form for Stacks material. No bibliography citation is involved.\par\medskip
\ifdefined\hypertarget\hypertarget{sourcepointer-2}{}\fi
\noindent\textit{Source reference.} Original lemma 010H: sources/stacks-homology/pinned-original-source.tex, source lines 1229--1229. Navigation pointer invoked at proof line 3445; the counted span contains the pointer only and the cited result is not reproduced inside it, so the pointer is registered as a non-load-bearing navigation stub in the established form for Stacks material. No bibliography citation is involved.\par\medskip
\ifdefined\hypertarget\hypertarget{sourcepointer-3}{}\fi
\noindent\textit{Source reference.} Original lemma 08N7: sources/stacks-homology/pinned-original-source.tex, source lines 1347--1347. Navigation pointer invoked at proof line 3464; the counted span contains the pointer only and the cited result is not reproduced inside it, so the pointer is registered as a non-load-bearing navigation stub in the established form for Stacks material. No bibliography citation is involved.\par\medskip
\begin{lemma}
\label{lemma-long-exact-sequence-cochain}
\begin{slogan}
Short exact sequences of complexes give rise to long exact sequences
of (co)homology.
\end{slogan}
Let $\mathcal{A}$ be an abelian category. Suppose that
$$
0 \to
A^\bullet \to
B^\bullet \to
C^\bullet \to
0
$$
is a short exact sequence of cochain complexes of $\mathcal{A}$.
Then there is a long exact cohomology sequence
$$
\xymatrix{
\ldots & \ldots & \ldots \ar[lld] \\
H^i(A^\bullet) \ar[r] &
H^i(B^\bullet) \ar[r] &
H^i(C^\bullet) \ar[lld] \\
H^{i + 1}(A^\bullet) \ar[r] &
H^{i + 1}(B^\bullet) \ar[r] &
H^{i + 1}(C^\bullet) \ar[lld] \\
\ldots & \ldots & \ldots \\
}
$$
The construction produces long exact cohomology sequences
which are functorial in the short exact
sequence and compatible with shifts as in
Definition \ref{definition-cohomology-shift}.
\end{lemma}

\begin{proof}
For the horizontal maps $H^i(A^\bullet) \to H^i(B^\bullet)$ and
$H^i(B^\bullet) \to H^i(C^\bullet)$ we use the fact that $H^i$
is a functor, see above. For the ``boundary map''
$H^i(C^\bullet) \to H^{i + 1}(A^\bullet)$ we use the map $\delta$
of the Snake Lemma \ref{lemma-snake}
applied to the diagram
$$
\xymatrix{
&
A^i/\Im(d_A^{i - 1}) \ar[r] \ar[d]^{d_A^i} &
B^i/\Im(d_B^{i - 1}) \ar[r] \ar[d]^{d_B^i} &
C^i/\Im(d_C^{i - 1}) \ar[r] \ar[d]^{d_C^i} &
0 \\
0 \ar[r] &
\Ker(d_A^{i + 1}) \ar[r] &
\Ker(d_B^{i + 1}) \ar[r] &
\Ker(d_C^{i + 1}) &
}
$$
This works as the kernel of the right vertical map is equal to
$H^i(C^\bullet)$ and the cokernel of the left vertical map is
$H^{i + 1}(A^\bullet)$. The exactness of the long sequence is
the exactnesss in part (2) of Lemma \ref{lemma-snake}.
The functoriality is Lemma \ref{lemma-snake-natural}. Compatibility
with shifts is immediate from the definitions.
\end{proof}

\end{document}
